\section{DIC Parameters}
\label{sec:parameters}

The \emph{Parameters} section of the left sidebar holds the subset-
and mesh-related choices. Change any of these and the downstream
pipeline will recompute when you press Run.

\subsection{Subset Size}

The IC-GN correlation window, in pixels. Displayed as an odd number
(e.g.\ 41), stored internally as the even half-width (40).

\begin{center}
\begin{tabular}{lp{8cm}}
  \toprule
  Value & When to use \\
  \midrule
  21--31 & Fine / dense speckle, small strain. Best spatial resolution. \\
  41 (default) & General speckle patterns with $\sim$ 3--5 px particles. \\
  51--81 & Noisy images, coarse speckle, or low-contrast regions. Lower spatial resolution. \\
  \bottomrule
\end{tabular}
\end{center}

Larger subsets average over more pixels, which improves signal-to-
noise but smooths out local features. Rule of thumb: each subset
should contain 5--10 speckle particles.

\subsection{Subset Step}

Node spacing in pixels, restricted to powers of two:
\texttt{4, 8, 16, 32, 64}. Smaller step means more mesh nodes and a
denser displacement field.

\begin{center}
\begin{tabular}{lp{8cm}}
  \toprule
  Value & Trade-off \\
  \midrule
  4 & Densest grid, very slow. Use only for $\leq$ 512\textsuperscript{2} images. \\
  8 (default for 512\textsuperscript{2}) & Good balance. \\
  16 (default for 1024\textsuperscript{2}) & Standard choice for large images. \\
  32--64 & Very fast, coarse. Only for quick previews. \\
  \bottomrule
\end{tabular}
\end{center}

\begin{tipbox}
Total mesh nodes grow as $(\text{ROI area}) / (\text{step})^2$.
Halving the step quadruples the node count and roughly quadruples
runtime. Start coarse, iterate on Region of Interest and parameters,
then refine the step only when you are happy with the setup.
\end{tipbox}

\subsection{Search Range}

Maximum per-frame displacement (in pixels) that the FFT initial
search can detect. If the real displacement exceeds this, the FFT
returns a clipped peak and pyALDIC auto-expands the search radius
($2\times$ each retry up to image half-size) --- but this costs
time.

\begin{center}
\begin{tabular}{lp{8cm}}
  \toprule
  Value & When to use \\
  \midrule
  8--20 (default 20) & Typical quasi-static experiments. \\
  40--60 & Larger translations, moderate rotations in incremental mode. \\
  80+ & Use only if you know the motion is huge. Diagnose before setting high: a large search range also means more chance of picking wrong NCC peaks in repetitive textures. \\
  \bottomrule
\end{tabular}
\end{center}

\begin{warnbox}
If you see \texttt{FFT search: N nodes have peaks at search boundary}
in the console log, the initial search range is too small. Either
raise it manually or let Auto-expand (in Advanced, on by default)
handle it.
\end{warnbox}

\subsection{Mesh Refinement}

Two independent checkboxes below the search range control quadtree
refinement of the mesh:

\begin{description}[leftmargin=1.2cm,style=nextline]
  \item[Refine Inner Boundary]
    Subdivide mesh elements near \emph{holes} inside the Region of
    Interest --- e.g.\ around a bubble, void, or cut-out region.
  \item[Refine Outer Boundary]
    Subdivide mesh elements near the outer edge of the Region of
    Interest, where boundary effects matter most.
\end{description}

When either box is checked (or the Refine brush has been used on
frame 1), the \emph{Refinement Level} dropdown activates:

\begin{center}
\begin{tabular}{lll}
  \toprule
  Level & Label & Min element size \\
  \midrule
  1 & Light        & $\text{step} / 2$ \\
  2 & Medium       & $\text{step} / 4$ \\
  3 & Heavy        & $\text{step} / 8$ \\
  4 & Extra Heavy  & $\text{step} / 16$ \\
  5 & Ultra        & $\text{step} / 32$ \\
  \bottomrule
\end{tabular}
\end{center}

Available levels depend on subset size and subset step so the min
element size stays reasonable (minimum 2 px element). A live readout
below the dropdown tells you the current minimum element size.
